\section{Appendix}

\subsection{Repository Layout}

Each primary loop skill follows the same layout:
\begin{verbatim}
<name>-loop/
  README.md
  SKILL.md
  prompt.md
  scripts/
    init_<name>_loop.cjs
    run_daemon.py
    run_daemon.sh
    start_<name>_loop.sh
    stop_<name>_loop.sh
    status_<name>_loop.sh
    print_cron_entry.sh
\end{verbatim}

\subsection{Example Active Task Schema}

\begin{verbatim}
{
  "plan_id": "PROJECT_ROADMAP",
  "mode": "optimize",
  "bounded_target": "complete one verifiable slice",
  "prompt_patch": [],
  "scope_patch": [],
  "verification_patch": [],
  "handoff": "next smallest useful action"
}
\end{verbatim}

\subsection{Example Failure Memory Entry}

\begin{verbatim}
{
  "failure_type": "verification_gap",
  "symptom": "implementation completed without running target tests",
  "patch": "before final handoff, run the smallest relevant check",
  "seen_in": ["epoch-004", "epoch-009"]
}
\end{verbatim}

\subsection{Daemon Management Commands}

All loop skills share the same management interface:

\begin{verbatim}
# Initialize state
node <skill>-loop/scripts/init_<name>_loop.cjs MY_ROADMAP

# Start daemon (tmux/nohup auto-select)
bash <skill>-loop/scripts/start_<name>_loop.sh MY_ROADMAP

# Check status (PID health + last mode)
bash <skill>-loop/scripts/status_<name>_loop.sh MY_ROADMAP

# Stop daemon gracefully (SIGTERM -> SIGKILL)
bash <skill>-loop/scripts/stop_<name>_loop.sh MY_ROADMAP

# Print cron entries for OS-level supervision
bash <skill>-loop/scripts/print_cron_entry.sh
\end{verbatim}

\subsection{Cron Integration Example}

The following crontab entries provide automatic restart after reboot and periodic health checks:
\begin{verbatim}
# Auto-start after reboot
@reboot cd /workspace && <SKILL>_LOOP_WORKSPACE=/workspace \
  bash /path/to/start_<name>_loop.sh

# Health check every N minutes (restart on failure)
*/N * * * * cd /workspace && \
  bash /path/to/status_<name>_loop.sh \
  || bash /path/to/start_<name>_loop.sh
\end{verbatim}

\subsection{Implementation Notes}

The Codex and MOA loops include additional monitoring commands:
\begin{verbatim}
# Codex: real-time monitoring dashboard
bash codex-loop/scripts/monitor_codex_loop.sh --watch

# MOA: DAG + Blackboard + PEFT dashboard
bash moa-loop/scripts/monitor_moa_loop.sh --watch
\end{verbatim}

The MOA loop extends the abstraction with \texttt{dag\_scheduler.py}, \texttt{shared\_blackboard.py}, and \texttt{iteration\_manager.py}. These modules implement horizontal DAG scheduling, centralized agent communication, and epoch-level PEFT adaptation tracking.
